% ==============================================================================
% Fichier : chapitres/03_securite_actifs.tex
% Livre III : Sécurité des Actifs
% ==============================================================================

\chapter{Sécurité des Actifs}
\label{chap:securite_actifs}

\begin{boxobjectifs}
\textbf{Objectifs du Livre~III :} Donner les éléments de compréhension liés aux actifs informationnels, leur importance et leur protection. À l'issue, l'étudiant saura classifier les données, appliquer les modèles de contrôle d'accès appropriés, gérer la rémanence et se référer aux standards internationaux de contrôle.
\end{boxobjectifs}

% ------------------------------------------------------------------------------
\section{Classification des Données}
% ------------------------------------------------------------------------------

\subsection{Niveaux de Sensibilité}

La classification des données permet d'appliquer des mesures de protection proportionnées à leur valeur et à leur sensibilité.

\begin{table}[H]
\centering
\begin{tabular}{p{3cm}p{5cm}p{4cm}}
    \toprule
    \textbf{Niveau} & \textbf{Description} & \textbf{Exemples} \\
    \midrule
    \textbf{Public}        & Aucune restriction de diffusion. & Site web institutionnel, communiqués de presse. \\
    \textbf{Interne}       & Usage réservé aux employés. & Procédures internes, organigrammes. \\
    \textbf{Confidentiel}  & Accès restreint aux personnes autorisées. & Contrats clients, données RH. \\
    \textbf{Secret/Très Sensible} & Accès ultra-restreint, impact majeur en cas de divulgation. & Données financières stratégiques, secrets industriels. \\
    \bottomrule
\end{tabular}
\caption{Niveaux de classification des données}
\end{table}

\subsection{Méthodes de Classification}
\begin{itemize}
    \item \textbf{Classification manuelle :} Un responsable désigne le niveau lors de la création.
    \item \textbf{Classification automatique :} Outils DLP (Data Loss Prevention) analysant le contenu et attribuant automatiquement un niveau.
    \item \textbf{Classification basée sur les rôles :} Le niveau dépend du rôle de l'utilisateur créateur.
\end{itemize}

% ------------------------------------------------------------------------------
\section{Protection des Données}
% ------------------------------------------------------------------------------

\subsection{Modèles de Contrôles d'Accès}

\begin{table}[H]
\centering
\small
\begin{tabular}{p{2.5cm}p{4.5cm}p{4.5cm}}
    \toprule
    \textbf{Modèle} & \textbf{Principe} & \textbf{Usage typique} \\
    \midrule
    \textbf{MAC} (Mandatory) & Niveau de sécurité attribué aux sujets et objets par l'administrateur système. L'accès dépend de la comparaison des étiquettes. & Systèmes gouvernementaux, militaires. \\
    \midrule
    \textbf{DAC} (Discretionary) & Le propriétaire de la ressource définit qui peut y accéder. & Systèmes de fichiers courants (Windows NTFS, Linux). \\
    \midrule
    \textbf{RBAC} (Role-Based) & Les droits sont attribués à des rôles ; les utilisateurs héritent des droits de leur rôle. & Applications d'entreprise (ERP, CRM). \\
    \midrule
    \textbf{ABAC} (Attribute-Based) & Accès basé sur les attributs du sujet, de l'objet et de l'environnement. & Cloud, IoT, systèmes distribués complexes. \\
    \bottomrule
\end{tabular}
\caption{Modèles de contrôle d'accès}
\end{table}

\subsection{Mesures Technico-logiques}
\begin{itemize}
    \item \textbf{Pare-feux :} Filtrage du trafic réseau.
    \item \textbf{Antivirus / EDR :} Détection et neutralisation des malwares.
    \item \textbf{Chiffrement :} Protection des données en transit (TLS) et au repos (AES).
    \item \textbf{DLP (Data Loss Prevention) :} Surveillance et blocage des fuites de données.
\end{itemize}

% ------------------------------------------------------------------------------
\section{Appartenance des Données (Data Ownership)}
% ------------------------------------------------------------------------------

\subsection{Définition et Importance}
L'appartenance des données désigne la propriété et la responsabilité de celles-ci au sein d'une organisation. Identifier clairement le propriétaire d'une donnée permet de :
\begin{itemize}
    \item Définir les rôles et responsabilités en matière de gestion.
    \item Établir des droits d'accès appropriés.
    \item Faciliter la traçabilité.
    \item Assurer la conformité réglementaire.
\end{itemize}

\subsection{Méthodes pour Déterminer l'Appartenance}
\begin{itemize}
    \item Identification du propriétaire (créateur ou responsable métier de la donnée).
    \item Analyse des flux de données au sein de l'organisation.
    \item Exploitation des métadonnées (auteur, date, département).
\end{itemize}

\subsection{Bonnes Pratiques}
\begin{itemize}
    \item Sensibiliser les employés à leurs responsabilités.
    \item Définir des politiques de classification et de gestion des accès claires.
    \item Réaliser des audits réguliers de la gestion des accès.
\end{itemize}

% ------------------------------------------------------------------------------
\section{Mémoire, Rémanence et Destruction}
% ------------------------------------------------------------------------------

\begin{boxdef}
\textbf{Rémanence des données :} Persistance de données dans un support de stockage, même après leur suppression apparente. Ce phénomène peut conduire à des fuites d'informations sensibles si les supports ne sont pas correctement effacés avant réaffectation ou destruction.
\end{boxdef}

\subsection{Risques associés à la Rémanence}
\begin{itemize}
    \item Vol de données lors de la revente ou du recyclage de matériel.
    \item Réutilisation inappropriée de supports sans effacement sécurisé.
    \item Non-conformité réglementaire (conservation de données au-delà de la durée légale).
\end{itemize}

\subsection{Méthodes d'Effacement Sécurisé}

\begin{table}[H]
\centering
\begin{tabular}{p{3cm}p{9cm}}
    \toprule
    \textbf{Méthode} & \textbf{Description} \\
    \midrule
    \textbf{Écrasement}   & Réécriture multiple sur les secteurs (normes DoD~5220.22-M : 3 passes minimum). \\
    \textbf{Dégaussement} & Exposition à un champ magnétique puissant détruisant les données magnétiques. \\
    \textbf{Destruction physique} & Déchiquetage, broyage ou incinération du support. Méthode la plus sûre. \\
    \textbf{Chiffrement + effacement} & Chiffrement préalable des données puis suppression de la clé. \\
    \bottomrule
\end{tabular}
\caption{Méthodes d'effacement sécurisé des données}
\end{table}

% ------------------------------------------------------------------------------
\section{Standards et Référentiels de Contrôle}
% ------------------------------------------------------------------------------

\begin{table}[H]
\centering
\small
\begin{tabular}{p{3.5cm}p{8.5cm}}
    \toprule
    \textbf{Standard} & \textbf{Description} \\
    \midrule
    \textbf{ISO/IEC~27001} & Système de Management de la Sécurité de l'Information (SMSI). Certification internationale de référence. \\
    \textbf{NIST SP~800-53} & Contrôles de sécurité pour les systèmes d'information fédéraux américains (18 familles de contrôles). \\
    \textbf{COBIT~2019}     & Cadre de gouvernance et de gestion des TI développé par l'ISACA (5 domaines). \\
    \textbf{CIS Controls}   & 18 contrôles de sécurité essentiels prioritaires, du Center for Internet Security. \\
    \textbf{PCI~DSS}        & Norme de sécurité des données pour l'industrie des cartes de paiement (12 exigences). \\
    \bottomrule
\end{tabular}
\caption{Principaux standards et référentiels de contrôle}
\end{table}

% ------------------------------------------------------------------------------
\section{Évaluation des Acquis --- Livre~III}
% ------------------------------------------------------------------------------

\begin{boxexo}
\textbf{QCM :}
\begin{enumerate}
    \item Quel modèle de contrôle d'accès est le plus adapté à un système militaire ? \textbf{(a) DAC \quad (b) RBAC \quad (c) MAC \quad (d) ABAC}
    \item La rémanence des données désigne : \textbf{(a) La durée de vie d'un disque dur \quad (b) La persistance de données après leur suppression apparente \quad (c) La sauvegarde automatique des données \quad (d) La compression des données}
    \item ISO/IEC 27001 porte sur : \textbf{(a) Les contrôles financiers \quad (b) Le système de management de la sécurité de l'information \quad (c) La gestion de projet \quad (d) La virtualisation}
    \item Quelle méthode d'effacement est la plus sûre pour un disque dur contenant des données classifiées Secret ? \textbf{(a) Formatage simple \quad (b) Suppression des fichiers \quad (c) Destruction physique \quad (d) Chiffrement ultérieur}
\end{enumerate}

\textbf{Travail Pratique :}\\
Proposer une politique de classification des données pour une université, en identifiant au moins cinq types de données traitées, leurs niveaux de sensibilité et les mesures de protection à appliquer.
\end{boxexo}
